\scanpage{25}\mspage{-- 23}
ne peut s'interpréter directement comme un chemin
\unsure{au|en} notre sens, il faut payer nos fautes ! Mais
on va à une telle flèche associer le chemin
\diagram{\begin{tikzcd}[ampersand replacement=\&,column sep=1.15em]
x \arrow[r] \& y \arrow[r,<->,"{\mathrm{id}}"] \& y \arrow[r,"{\mathrm{id}}"] \& y \arrow[r,<->,"{\mathrm{id}}"] \& y
\end{tikzcd}},
\deleted{de} \inserted{élément de} $\underline{\unsure{\mathrm{Ch}|\mathscr{L}}}_{\unsure{\Sigma|Z},2}(x,y)$, ce qui définit l'application \unsure{cherchée|du début}.
Il faut vérifier la compatibilité avec les identités
et avec la composition. À la flèche $1_x$ \inserted{$\colon x\mathord{-}x$} est associé
l'élément de $\unsure{\mathscr{L}|\mathcal{L}}_2(x,x)$
\diagram{\begin{tikzcd}[ampersand replacement=\&,column sep=1.15em]
x \arrow[r,no head,"{\mathrm{id}}"] \& x \arrow[r,<-,"{\mathrm{id}}"] \& x \arrow[r,"{\mathrm{id}}"] \& x \arrow[r,<-,"{\mathrm{id}}"] \& x
\end{tikzcd}},
Il faut prouver que dans $\pi_0(\underline{\mathrm{Ch}}_{\Sigma,2}(x,x))$ il devient
égal à l'élément déduit de $1_x\in\mathscr{L}_{\Sigma,0}(x,x)$,
par le foncteur de transition. Mais cet élément
n'est autre justement que
\diagram{\begin{tikzcd}[ampersand replacement=\&,column sep=1.15em]
x \arrow[r,"{\mathrm{id}}"] \& x \arrow[r,<-,"{\mathrm{id}}"] \& x \arrow[r,"{\mathrm{id}}"] \& x \arrow[r,<-,"{\mathrm{id}}"] \& x
\end{tikzcd}},
c'est le $=$ ! (sans avoir $=$ à passer au $\pi_0$) !

Pour la compatibilité avec la composition
\diagram{\begin{tikzcd}[ampersand replacement=\&,column sep=1.15em]
x \arrow[r,"{u}"] \& y \arrow[r,no head,"{v}"] \& z \arrow[r,<-] \& {}
\end{tikzcd}}
On a d'une part
\diagram{\begin{tikzcd}[ampersand replacement=\&,column sep=1.15em]
x \arrow[r,no head,"{vu}"] \& z \arrow[r,<-,"{\mathrm{id}}"] \& z \arrow[r,no head,"{\mathrm{id}}"] \& z \arrow[r,<-,"{\mathrm{id}}"] \& z
\end{tikzcd}}
d'\deleted{$x\to y\mathord{-}$}autre part
\diagram{\begin{tikzcd}[ampersand replacement=\&,column sep=1.15em]
x \arrow[r,"{u}"] \& y \arrow[r,<-,"{\mathrm{id}}"] \& y \arrow[r,"{\mathrm{id}}"] \& y \arrow[r,<-,"{\mathrm{id}}"] \& y \arrow[r,"{v}"] \& z \arrow[r,<-,"{\mathrm{id}}"] \& z \arrow[r,"{\mathrm{id}}"] \& z \arrow[r,<-,"{\mathrm{id}}"] \& z
\end{tikzcd}}
on complète le premier chemin par dégénérescence
aux deux bouts, pour le comparer au deuxième,
en tant qu'éléments de $\underline{\unsure{\mathrm{Ch}|\mathscr{L}}}_4(x,z)$ ;
\diagram{\begin{tikzcd}[ampersand replacement=\&,column sep=1.15em]
x \arrow[r,"{\mathrm{id}}"] \& x \arrow[r,<-,"{\mathrm{id}}"] \& x \arrow[r,no head,"{vu}"] \& z \arrow[r,<->,"{\mathrm{id}}"] \& z \arrow[r,<->,"{\mathrm{id}}"] \& z \arrow[r,<-,"{\mathrm{id}}"] \& z \arrow[r,"{\mathrm{id}}"] \& z \arrow[r,<-,"{\mathrm{id}}"] \& z
\end{tikzcd}}
% archive sheet number: 24
